%======================================================================
\section{Verificação e plano de testes}\label{sec:verificacao}
%======================================================================
\subsection{Níveis de verificação}

A verificação do sistema é feita em três níveis (Tabela~\ref{tab:niveis}), dos testes
automáticos no computador até os testes na placa, com os sensores.

\begin{table}[H]
\centering
\caption{Níveis de verificação.}
\label{tab:niveis}
\small
\begin{tabular}{@{}l>{\raggedright\arraybackslash}p{11.0cm}@{}}
\toprule
Nível & Significa \\
\midrule
testes no PC & testes automáticos rodando no computador (lógica pura, sem placa) \\
ponta a ponta sem placa & central real em Docker, com o simulador no lugar do hub \\
bancada & rodando na placa, com os sensores \\
\bottomrule
\end{tabular}
\end{table}

\subsection{Verificações realizadas}

\begin{table}[H]
\centering
\caption{Como cada parte foi verificada.}
\label{tab:verificado}
\small
\begin{tabular}{@{}>{\raggedright\arraybackslash}p{4.0cm}>{\raggedright\arraybackslash}p{9.8cm}@{}}
\toprule
Parte & Como \\
\midrule
Firmware do hub & testes no computador das regras de alerta (82 verificações), da configuração por MQTT (60) e do envio por mudança (23) \\
Firmware da pulseira & testes no computador do pipeline de queda (885 verificações, com a paridade C $\times$ Python), do processamento de sinal (25) e dos gestos do botão (42) \\
Mapa de pinos, chaves e MAC & pino repetido ou proibido e MAC mal escrito fazem o build falhar \\
Modelo de queda & validação por pessoa na SisFall; paridade C $\times$ Python \\
Central & os 5 serviços sobem; o broker recusa conexão anônima e senha errada; o simulador publicou estados e alertas, o InfluxDB gravou, as consultas do painel rodaram e o notificador formatou cada episódio \\
Notificador e demonstração & 22 e 19 testes automáticos \\
PCB v1 & \ercViolacoes{} violações de ERC e \drcErros{} erros de DRC; encaixe e ordem dos pinos no teste no papel 1:1 \\
Pulseira na placa & perfil de demonstração: 200\,Hz contínuos por 16\,min e $\pm$8\,g confirmado (abaixo) \\
\bottomrule
\end{tabular}
\end{table}

\paragraph{Pulseira na placa.} Em 13/11/2026, a pulseira com o perfil de demonstração,
ligado ao computador por USB, mandou as amostras brutas a 200\,Hz por 16 minutos sem parar.
Parado, $\amag$ ficou em $\approx$4077 contagens, ou seja, $\approx$1\,g na escala de
4096 contagens por g: a faixa de $\pm$8\,g está configurada. No mesmo ensaio, o programa de
demonstração parou de atualizar quando o link USB piscou, e a leitura da serial foi corrigida
para reabrir a porta sozinha.

\subsection{Plano de testes em bancada}

A Tabela~\ref{tab:plano} lista os testes na placa, com os sensores, na ordem em que fazem
sentido, com o critério de aprovação. As metas vêm dos requisitos (Seção~\ref{sec:requisitos}).

\begin{table}[H]
\centering
\caption{Plano de testes em bancada.}
\label{tab:plano}
\small
\begin{tabular}{@{}>{\raggedright\arraybackslash}p{2.7cm}>{\raggedright\arraybackslash}p{6.0cm}>{\raggedright\arraybackslash}p{5.1cm}@{}}
\toprule
Teste & Procedimento & Critério \\
\midrule
Amostragem & modo de coleta; intervalos entre amostras & 200\,Hz, \textit{jitter} $\le \pm$5\,ms (RNF01) \\
Faixa do MPU & placa parada; uma batida forte & $\sim$4096 contagens no eixo vertical; batida acima de 3,46\,g sem teto plano \\
Pipeline na placa & balançar; soltar a placa numa almofada; comparar o CSV bruto com o Python & disparos com P baixa; queda detectada; mesmos picos e mesma P do computador \\
Portas & abrir e fechar depressa, várias vezes & nenhuma transição perdida ou inventada \\
PIR & movimento depois dos 60\,s de aquecimento & detecção e volta a ``livre'' \\
Temperatura & 3 ou 4 pontos (perto de 10, 20 e 30\,°C) contra um termômetro; depois, termopar; 30\,min com o quarto parado & $\pm$1\,°C (RNF10); parado, o ruído fica abaixo de 0,2\,°C e sai uma mensagem a cada 5\,min \\
Peças ausentes & sem SSID, senha errada, broker parado, LM35 solto, MPU solto, CAN sem transceptor & um aviso por caso, e o resto funcionando \\
CAN entre as placas & pulseira com CAN de bancada, transceptores e terminação; tirar um fio do barramento & quadros \code{0x0FD}/\code{0x0FE} no log do hub; um aviso ao cair e um ao voltar \\
Botão & aperto simples; dois toques sem alerta; dois toques na janela de uma queda; segurar 3\,s; segurar durante um alerta; ficar com o botão apertado por 30\,s & só os dois toques na janela cancelam e só o segurar fora de alerta pede ajuda; o resto não faz nada; segurar sem soltar dá um pedido só \\
ESP-NOW & boot da pulseira; queda simulada (dois toques, opção de teste); ``estou bem''; pedido de ajuda & canal descoberto; heartbeat no Grafana; alertas \code{queda} e \code{ajuda}, \code{ativo} e \code{cancelado} \\
Queda atrasada & hub desligado, queda, hub religado em $\sim$20\,s & o alerta chega uma vez só \\
Queda após reinício & queda, pulseira reiniciada, nova queda & dois alertas, com \code{id} diferentes \\
Remetente estranho & outro ESP32 mandando ESP-NOW ao hub & um aviso e nenhum alerta \\
Latências & hora na pulseira, no hub, no broker e no celular & pico $\to$ MQTT $\le$ 500\,ms; porta $\to$ MQTT $\le$ 200\,ms; alerta $\to$ Telegram $\le$ 5\,s \\
\textit{Last Will} & desligar o hub & \code{status} \code{offline} em até $\sim$45\,s \\
Notificador & derrubar a internet 1\,min com um alerta no meio; reiniciar o contêiner & os dois alertas chegam \\
Disponibilidade & 24\,h contínuas, derrubando a rede & reconexão sozinha, nenhum evento QoS 1 perdido \\
Autonomia & medidor USB entre o \textit{power bank} e a pulseira & consumo medido (RNF14) \\
\bottomrule
\end{tabular}
\end{table}
